\section{Introduction}
\label{sec:introduction}

Aerial 3D scene understanding is central to autonomous uncrewed aerial vehicle (UAV) operation in cargo delivery, search and rescue, inspection, and human-UAV collaboration. In these settings, a perception system must infer not only visible image semantics, but also metric depth, free space, and the semantic occupancy of partially observed volumes. Semantic occupancy prediction combines partial visual observations with learned scene priors to infer the spatial structure and semantic composition of unobserved regions, supporting anticipatory action planning under uncertainty. These plans can be refined in real time as additional visual evidence becomes available. RGB cameras are attractive for UAVs because they are lightweight, inexpensive, and compatible with strict size, weight, and power constraints. However, monocular 3D perception remains difficult: image formation collapses scale, depth, and occlusion into a two-dimensional signal, and aerial views introduce altitude-dependent scale changes, oblique/top-down viewpoints, weak texture, and occlusion patterns that differ from ground-level driving scenes.

Most progress in semantic scene completion (SSC) and camera-based 3D perception has been driven by ground-level indoor and autonomous-driving benchmarks. SSCNet introduced semantic completion from depth input \cite{song2017sscnet}; MonoScene extended the task to monocular RGB \cite{cao2022monoscene}; and subsequent methods such as VoxFormer \cite{li2023voxformer}, TPVFormer \cite{huang2023tpvformer}, SurroundOcc \cite{wei2023surroundocc}, OpenOccupancy \cite{wang2023openoccupancy}, CGFormer \cite{yu2024cgformer}, and FoundationSSC \cite{chen2025foundationssc} improved view transformation, voxel attention, and semantic--geometric fusion. More recently, VoxDet reformulated voxel-wise SSC as dense object detection by decoupling semantic prediction from offset regression. Its training-free Voxel-to-Instance (VoxNT) transformation converts voxel-level semantic labels into six-direction object-boundary offset targets; the predicted offsets then guide instance-aware feature aggregation in the semantic branch \cite{li2025voxdet}. Together, these designs demonstrate the value of explicit 3D representations while showing that completion can benefit from both geometry-aware lifting and instance-aware voxel reasoning.

Existing aerial datasets provide valuable but complementary forms of supervision. UAVid focuses on 2D semantic segmentation of oblique urban imagery \cite{lyu2020uavid}; Mid-Air and SynDrone provide synthetic multimodal flights with depth and semantic labels \cite{fonder2019midair,rizzoli2023syndrone}; Dronescapes supports semantic segmentation, metric depth, and surface-normal estimation from real UAV video \cite{marcu2023dronescapes}; UrbanScene3D targets aerial path planning and multiview reconstruction \cite{lin2022urbanscene3d}; and UAVScenes combines image and LiDAR semantics with calibrated poses for several 2D and 3D perception tasks \cite{wang2025uavscenes}. These resources can benchmark segmentation, MDE, or reconstruction, but they do not directly provide the task-aligned dense semantic occupancy and free-space voxel ground truth required for SSC evaluation. OccuFly \cite{gross2026occufly} fills this gap by jointly providing real aerial RGB observations, metric depth maps, camera calibration and poses, altitude metadata, and semantic occupancy grids across multiple scenes and flight heights. We therefore use OccuFly to evaluate MDE, calibrated multiview geometry, and SSC under a common aerial protocol, asking whether modern foundation representations supply transferable geometric and semantic priors and whether neighboring views improve SSC beyond monocular depth and visual-context priors.

We study the transfer of frozen V-JEPA~2.1 and DINOv3 representations to aerial monocular depth estimation (MDE) and depth-aware semantic scene completion (SSC). Whereas V-JEPA learns by predicting latent representations rather than reconstructing pixels \cite{bardes2024vjepa,assran2025vjepa2}, V-JEPA~2.1 augments this objective with dense prediction over both masked and visible tokens and with deep self-supervision across encoder levels. The latter propagates local spatial information to the output representation, enabling the final encoder layer alone to support dense downstream prediction \cite[Secs.~2.3.1--2.3.2 and App.~9.1]{murlabadia2026vjepa21}. DINOv3 provides a complementary image representation optimized for dense visual transfer \cite{simeoni2025dinov3}. Using these frozen backbones, we examine how linearly accessible depth, multiscale decoding \cite{ranftl2021dpt,li2022vitdet}, camera-aware feature modulation \cite{perez2018film}, and temporal aggregation affect aerial MDE. We then assess how monocular and spatio-temporal visual priors, calibrated multiview geometry from a frozen MVSFormer++ depth estimator \cite{cao2024mvsformerpp}, and instance-aware voxel reasoning \cite{li2025voxdet} contribute to SSC. Only task-specific prediction and fusion modules are trained. Controlled comparisons isolate individual design factors, while repeated independent training runs quantify variability for the principal model; results for architectural controls evaluated without repetition are interpreted descriptively rather than causally.

Our evaluation has three parts. First, we compare linear probes, DPT-style multiscale decoders \cite{ranftl2021dpt}, Simple Feature Pyramid (SFP) adapters inspired by ViTDet \cite{li2022vitdet}, and Feature-wise Linear Modulation (FiLM) \cite{perez2018film} using altitude and camera intrinsics. This separates linear depth accessibility, dense reconstruction capacity, final-layer pyramid adaptation, and metric scale conditioning. Second, we compare RGB features with spatio-temporal V-JEPA~2.1 video features using a dense temporal-attention adapter that converts tubelet tokens into target-frame maps. Third, we evaluate SSC variants inspired by CGFormer and FoundationSSC, replacing or augmenting their geometry pathways with V-JEPA-derived depth probabilities and a pose-aware multiview stereo (MVS) cost-volume branch based on calibrated UAV views.

Our contributions are:
\begin{itemize}
    \item A systematic frozen-backbone evaluation of V-JEPA~2.1 RGB and video representations for aerial metric depth estimation on OccuFly.
    \item A controlled comparison of linear probes, DPT-style decoders, SFP-adapted final-layer decoders, and altitude/intrinsics-conditioned FiLM variants.
    \item An analysis of spatio-temporal feature transfer through dense temporal attention over V-JEPA~2.1 video tokens.
    \item A depth-to-SSC study that injects V-JEPA-derived depth maps and depth distributions into CGFormer- and FoundationSSC-inspired aerial SSC variants.
    \item A pose-aware FoundationSSC variant that replaces monocular depth lifting with calibrated multiview cost-volume depth estimation from neighboring UAV frames.
\end{itemize}
